\subsection{Salary/Payroll টেবিলে বেতন হিসাব: রাশি ও সারসংক্ষেপের Query}

নিচের সারণিতে HR (Human Resources) মানবসম্পদ বিভাগ এবং IT (Information Technology) তথ্যপ্রযুক্তি বিভাগ বোঝায়।

\begin{flushleft}\begin{minipage}{\linewidth}\raggedright

\scriptsize

\begin{tabular}{|l|l|l|l|l|l|l|}
\hline

\chSixHeader
EmpID & Name & Department & BasicSalary & Allowance & Medical & Deduction\\
\hline
E101 & Rafi & Accounts & 30000 & 5000 & 1500 & 2000\\\hline
E102 & Nila & HR & 35000 & 6000 & 1800 & 2500\\\hline
E103 & Sami & IT & 45000 & 8000 & 2200 & 3000\\\hline
E104 & Mitu & Accounts & 28000 & 4500 & 1400 & 1500\\\hline
E105 & Hasan & IT & 50000 & 9000 & 2500 & 3500\\
\hline

\end{tabular}

\captionof{table}{Salary Table}

\end{minipage}\end{flushleft}

\chSixExample{৭.৫৮}{প্রত্যেক কর্মীর প্রাপ্য বেতন বা Net Salary দেখাও।}

ধরি—

\par NetSalary = BasicSalary + Allowance + Medical - Deduction.\par

SQL—

\noindent\begin{minipage}{\linewidth}
\noindent\begin{minipage}[c]{.53\linewidth}
\begin{verbatim}
SELECT Name, BasicSalary + Allowance
       + Medical - Deduction AS NetSalary
FROM Salary;
\end{verbatim}
\end{minipage}\hfill\begin{minipage}[c]{.04\linewidth}\raggedright{\color{chSixTeal}$\Rightarrow$}\end{minipage}\hfill
\begin{minipage}[c]{.40\linewidth}\raggedright\footnotesize\setlength{\tabcolsep}{2pt}\textbf{ফল}\par\smallskip
\chSixResultList{Name, NetSalary}{\mbox{Rafi, 34500}; \mbox{Nila, 40300}; \mbox{Sami, 52200}; \mbox{Mitu, 32400}; \mbox{Hasan, 58000}}
\end{minipage}
\end{minipage}\par\smallskip

এখানে NetSalary কুয়েরির মাধ্যমে হিসাব করা ফল। মূল টেবিলে এই নামে আলাদা ফিল্ড থাকা বাধ্যতামূলক নয়।

\noindent\begin{minipage}{\linewidth}
\chSixExample{৭.৫৯}{যেসব কর্মীর Basic Salary 35000-এর বেশি, তাঁদের তথ্য দেখাও।}

\noindent\begin{minipage}[c]{.53\linewidth}
\begin{verbatim}
SELECT EmpID, Name, BasicSalary
FROM Salary
WHERE BasicSalary>35000;
\end{verbatim}
\end{minipage}\hfill\begin{minipage}[c]{.04\linewidth}\raggedright{\color{chSixTeal}$\Rightarrow$}\end{minipage}\hfill
\begin{minipage}[c]{.40\linewidth}\raggedright\footnotesize\setlength{\tabcolsep}{2pt}\textbf{ফল}\par\smallskip
\begin{tabular}{|l|l|l|}\hline\chSixHeader
EmpID & Name & BasicSalary\\\hline
E103 & Sami & 45000\\\hline
E105 & Hasan & 50000\\
\hline\end{tabular}
\end{minipage}
\end{minipage}\par\smallskip

ফলাফলে Sami ও Hasan-এর তথ্য পাওয়া যাবে।

\noindent\begin{minipage}{\linewidth}
\chSixExample{৭.৬০}{IT বিভাগের কর্মীদের তথ্য দেখাও।}

\noindent\begin{minipage}[c]{.53\linewidth}
\begin{verbatim}
SELECT Name, BasicSalary
FROM Salary
WHERE Department='IT';
\end{verbatim}
\end{minipage}\hfill\begin{minipage}[c]{.04\linewidth}\raggedright{\color{chSixTeal}$\Rightarrow$}\end{minipage}\hfill
\begin{minipage}[c]{.40\linewidth}\raggedright\footnotesize\setlength{\tabcolsep}{2pt}\textbf{ফল}\par\smallskip
\begin{tabular}{|l|l|}\hline\chSixHeader
Name & BasicSalary\\\hline
Sami & 45000\\\hline
Hasan & 50000\\
\hline\end{tabular}
\end{minipage}
\end{minipage}\par\smallskip

\noindent\begin{minipage}{\linewidth}
\chSixExample{৭.৬১}{Basic Salary-এর সর্বোচ্চ মান বের করো।}

\noindent\begin{minipage}[c]{.53\linewidth}
\begin{verbatim}
SELECT MAX(BasicSalary) AS HighestSalary
FROM Salary;
\end{verbatim}
\end{minipage}\hfill\begin{minipage}[c]{.04\linewidth}\raggedright{\color{chSixTeal}$\Rightarrow$}\end{minipage}\hfill
\begin{minipage}[c]{.40\linewidth}\raggedright\footnotesize\setlength{\tabcolsep}{2pt}\textbf{ফল}\par\smallskip
\begin{tabular}{|l|}\hline\chSixHeader
HighestSalary\\\hline
50000\\
\hline\end{tabular}
\end{minipage}
\end{minipage}\par\smallskip

\noindent\begin{minipage}{\linewidth}
\chSixExample{৭.৬২}{Basic Salary-এর গড় বের করো।}

\noindent\begin{minipage}[c]{.53\linewidth}
\begin{verbatim}
SELECT AVG(BasicSalary) AS AverageSalary
FROM Salary;
\end{verbatim}
\end{minipage}\hfill\begin{minipage}[c]{.04\linewidth}\raggedright{\color{chSixTeal}$\Rightarrow$}\end{minipage}\hfill
\begin{minipage}[c]{.40\linewidth}\raggedright\footnotesize\setlength{\tabcolsep}{2pt}\textbf{ফল}\par\smallskip
\begin{tabular}{|l|}\hline\chSixHeader
AverageSalary\\\hline
37600\\
\hline\end{tabular}
\end{minipage}
\end{minipage}\par\smallskip

\noindent\begin{minipage}{\linewidth}
\chSixExample{৭.৬৩}{প্রতিটি বিভাগের গড় Basic Salary দেখাও।}

\noindent\begin{minipage}[c]{.53\linewidth}
\begin{verbatim}
SELECT Department,
       AVG(BasicSalary) AS AverageSalary
FROM Salary
GROUP BY Department;
\end{verbatim}
\end{minipage}\hfill\begin{minipage}[c]{.04\linewidth}\raggedright{\color{chSixTeal}$\Rightarrow$}\end{minipage}\hfill
\begin{minipage}[c]{.40\linewidth}\raggedright\footnotesize\setlength{\tabcolsep}{2pt}\textbf{ফল}\par\smallskip
\begin{tabular}{|l|l|}\hline\chSixHeader
Department & AverageSalary\\\hline
Accounts & 29000\\\hline
HR & 35000\\\hline
IT & 47500\\
\hline\end{tabular}
\end{minipage}
\end{minipage}\par\smallskip

\noindent\begin{minipage}{\linewidth}
\chSixExample{৭.৬৪}{যেসব বিভাগের গড় Basic Salary 30000-এর বেশি, সেগুলো দেখাও।}

\noindent\begin{minipage}[c]{.53\linewidth}
\begin{verbatim}
SELECT Department,
       AVG(BasicSalary) AS AverageSalary
FROM Salary
GROUP BY Department
HAVING AVG(BasicSalary)>30000;
\end{verbatim}
\end{minipage}\hfill\begin{minipage}[c]{.04\linewidth}\raggedright{\color{chSixTeal}$\Rightarrow$}\end{minipage}\hfill
\begin{minipage}[c]{.40\linewidth}\raggedright\footnotesize\setlength{\tabcolsep}{2pt}\textbf{ফল}\par\smallskip
\begin{tabular}{|l|l|}\hline\chSixHeader
Department & AverageSalary\\\hline
HR & 35000\\\hline
IT & 47500\\
\hline\end{tabular}
\end{minipage}
\end{minipage}\par\smallskip

HR ও IT বিভাগের গড় বেতন শর্ত পূরণ করে, তাই এই বিভাগ দুটির তথ্য পাওয়া যাবে।

\Needspace{18\baselineskip}
